% 原文第36页；原式与电路图保留。
\zhbox{54}{64}{235}{25}{9.5}{电容 $C_{OUT}$ 应能减小输出电压纹波 $\Delta V_{OUT}$。}
\zhbox{54}{95}{235}{102}{9.5}{为降低 $\Delta V_{OUT}$ 并改善瞬态响应，可能需要更大的 $C_{OUT}$。建议采用低 ESR 陶瓷电容，其体积小、输出电压纹波低。对于陶瓷电容，在开关频率 $f_{SW}$ 处，阻抗主要由电容量决定，并由此产生大部分输出纹波。为简化计算，$\Delta V_{OUT}$ 可由式（14）估算：}
\zhbox{54}{260}{235}{51}{9.5}{对于混合型、聚合物或电解电容，在 $f_{SW}$ 处 ESR 在阻抗中占主导。为简化计算，$\Delta V_{OUT}$ 可由式（15）估算：}
\zhbox{54}{365}{235}{38}{9.5}{对于 100W USB PD 应用，建议使用一个 100$\mu$F 聚合物电容和四个 22$\mu$F 陶瓷电容。}
\zhbox{54}{406}{235}{68}{9.5}{选择满足设计中输出电压纹波及负载瞬态要求的 $C_{OUT}$。设计高输出电压应用时，应考虑电容量降额。}
\zhbox{54}{478}{235}{106}{9.5}{\textbf{外部 MOSFET 选择}\\MP2984 需要四个外部 N 沟道功率 MOSFET（见图12）。SWA 和 SWD 为两个上管，SWB 和 SWC 为两个下管。降压模式下，SWA 和 SWB 进行开关动作，SWD 导通。升压模式下，SWC 和 SWD 进行开关动作，SWA 导通。}
\zhbox{92}{708}{170}{13}{9.5}{\centering\textbf{图12：升降压拓扑}}
\zhbox{324}{64}{235}{26}{9.5}{选择 MOSFET 的关键参数如下：}
\zhbox{324}{96}{235}{102}{9.5}{\textbullet\ \textbf{最大漏源电压（$V_{DS\_MAX}$）：}SWA 和 SWB 必须承受最大输入电压及 SW1 处的瞬态尖峰。建议其 $V_{DS\_MAX}$ 取 $1.5\times V_{IN}$。SWC 和 SWD 必须承受输出电压及 SW2 处的瞬态尖峰，建议其 $V_{DS\_MAX}$ 不小于 $1.5\times V_{OUT}$。}
\zhbox{324}{201}{235}{64}{9.5}{\textbullet\ \textbf{最大连续漏极电流（$I_{D\_MAX}$）：}SWA、SWB、SWC 和 SWD 在开关过程中必须承受最大电感电流 $I_{L\_MAX}$ 和峰值电感电流 $I_{L\_PEAK}$。}
\zhbox{324}{268}{235}{78}{9.5}{\textbullet\ \textbf{阈值电压（$V_{TH}$）：}栅极驱动电压由 VCC 提供。MOSFET 的栅极平台电压不应超过最小 $V_{CC}$，否则在启动或过载（OL）条件下，MOSFET 可能无法充分导通。}
\zhbox{324}{348}{235}{51}{9.5}{\textbullet\ \textbf{总栅极电荷（$Q_G$）：}在栅极电压 7.2V 条件下，各 MOSFET 的 $Q_G$ 均应小于 50nC。若两个 MOSFET 并联，则每个的 $Q_G$ 应小于 25nC。}
\zhbox{324}{402}{235}{64}{9.5}{\textbullet\ \textbf{漏源导通电阻（$R_{DS(ON)}$）：}较低的导通电阻有利于控制 MOSFET 温升并提高效率。有关导通电阻的选择，请参阅 MOSFET 说明及式（16）至（24）。}
\zhbox{324}{471}{235}{54}{9.5}{\textbf{SWA MOSFET}\\升压模式下，SWA 导通，其导通功率损耗 $P_{CON\_LOSS\_SWA}$ 可由式（16）计算：}
\zhbox{324}{574}{235}{77}{9.5}{若 MOSFET 的结到环境热阻为 50°C/W（由电路板散热能力决定），最大允许温升为 50°C，则最大功率损耗 $P_{CON\_LOSS}$ 为 1W。$P_{CON\_LOSS\_SWA}$ 可由式（17）计算：}
\zhbox{324}{684}{235}{27}{9.5}{降压模式下，$P_{CON\_LOSS\_SWA}$ 可由式（18）计算：}

